\subsection{网络费用流}

\subsubsection{基于 EK 的 SSP 算法}

每次用 SPFA 找最短路（按费用）增广，基于 \texttt{EdmondsKarp}。复杂度上界 $\mathcal{O}(f\cdot nm)$，$f$ 为最大流值。

\begin{lstlisting}
using i64 = long long;
using p64 = std::array<i64,2>;
constexpr i64 INF = 1e18;
struct SSP {
    int n;
    struct Edge {
        int to = 0,nxt = 0;
        i64 f = 0,c = 0,w = 0;
    };
    std::vector<int> head,fr;
    std::vector<Edge> edge;
    std::vector<i64> dis;
    std::vector<bool> vis;
    SSP(int _n):n(_n),head(n+1),edge(2),fr(n+1),dis(n+1),vis(n+1){}
    void addEdge(int u,int v,i64 f,i64 c,i64 w){
        edge.push_back({v,head[u],f,c,w});
        head[u] = edge.size() - 1;
    }
    void connect(int u,int v,i64 c,i64 w){
        addEdge(u,v,c,c,w);
        addEdge(v,u,0,c,-w);
    }
    p64 work(int s,int t){
        p64 ans = {0,0};
        while (1) {
            std::queue<int> que;
            que.push(s);
            std::fill(fr.begin(),fr.end(),0);
            std::fill(vis.begin(),vis.end(),0);
            std::fill(dis.begin(),dis.end(),INF);
            dis[s] = 0,vis[s] = 1;
            while (!que.empty()) {
                int u = que.front();
                que.pop();
                vis[u] = 0;
                for (int e = head[u];e;e = edge[e].nxt){
                    i64 f = edge[e].f,w = edge[e].w;
                    int v = edge[e].to;
                    if (f > 0 && dis[u] + w < dis[v]){
                        dis[v] = dis[u] + w;
                        fr[v] = e^1;
                        if (!vis[v]){
                            que.push(v);
                            vis[v] = 1;
                        }
                    }
                }
            }
            if (dis[t] >= INF) break;
            i64 mn = INF;
            for (int u = t;u != s;u = edge[fr[u]].to){
                mn = std::min(mn,edge[fr[u] ^ 1].f);
            }
            for (int u = t;u != s;u = edge[fr[u]].to){
                edge[fr[u]].f += mn;
                edge[fr[u] ^ 1].f -= mn;
            }
            ans[0] += mn;
            ans[1] += dis[t] * mn;
        }
        return ans;
    }
};
\end{lstlisting}

\subsubsection{基于 Dinic 的 SSP 算法（Dijkstra 势能优化）}

先用 SPFA 求初始势能 $h$，再用势能将费用化为非负后跑 Dijkstra 找增广路，基于 \texttt{Dinic}。复杂度上界 $\mathcal{O}(f\cdot m\log n)$，$f$ 为最大流值。

注意：以上两种算法均无法处理带负环费用的最小费用最大流。

\begin{lstlisting}
using i64 = long long;
using p64 = std::array<i64,2>;
constexpr i64 INF = 1e18;
struct MinCostFlow {
    int n;
    struct Edge {
        int to = 0,nxt = 0;
        i64 f = 0,c = 0,w = 0;
    };
    std::vector<int> head,now;
    std::vector<Edge> edge;
    std::vector<i64> dis,h;
    std::vector<bool> vis;
    MinCostFlow (int _n):n(_n),head(n+1),edge(2),now(n+1),dis(n+1),h(n+1),vis(n+1){}
    void addEdge(int u,int v,i64 f,i64 c,i64 w){
        edge.push_back({v,head[u],f,c,w});
        head[u] = edge.size() - 1;
    }
    void connect(int u,int v,i64 c,i64 w){
        addEdge(u,v,c,c,w);
        addEdge(v,u,0,c,-w);
    }
    void SPFA(int s,int t){
        std::fill(h.begin(),h.end(),INF);
        std::fill(vis.begin(),vis.end(),0);
        std::queue<int> que;
        que.push(s);
        h[s] = 0,vis[s] = 1;
        while (!que.empty()){
            int u = que.front();
            que.pop();
            vis[u] = 0;
            for (int e = head[u];e;e = edge[e].nxt){
                i64 f = edge[e].f,w = edge[e].w;
                int v = edge[e].to;
                if (f > 0 && h[u] + w < h[v]){
                    h[v] = h[u] + w;
                    if (!vis[v]) {
                        que.push(v);
                        vis[v] = 1;
                    }
                }
            }
        }
    }
    bool dijkstra(int s,int t){
        std::fill(dis.begin(),dis.end(),INF);
        std::fill(vis.begin(),vis.end(),0);
        std::priority_queue<p64> pq;
        pq.push({0,s});
        dis[s] = 0;
        while (!pq.empty()){
            auto [dist,u] = pq.top();
            pq.pop();
            if (vis[u]) continue;
            vis[u] = 1;
            for (int e = head[u];e;e = edge[e].nxt){
                int v = edge[e].to;
                i64 w = edge[e].w,f = edge[e].f;
                if (f > 0 && dis[u] + w + h[u] - h[v] < dis[v]){
                    dis[v] = dis[u] + w + h[u] - h[v];
                    pq.push({-dis[v],v});
                }
            }
        }
        return vis[t];
    }
    i64 dfs(int u,const int& t,i64 flow,i64 &cost){
        if (u == t) return flow;
        i64 sum = 0;
        vis[u] = 1;
        for (int e = now[u];e;e = edge[e].nxt){
            now[u] = e;
            int v = edge[e].to;
            i64 f = edge[e].f,w = edge[e].w;
            if (f > 0 && !vis[v] && h[v] == h[u] + w) {
                i64 f2 = dfs(v,t,std::min(f,flow),cost);
                edge[e].f -= f2;
                edge[e^1].f += f2;
                sum += f2;
                flow -= f2;
                cost += f2 * w;
                if (flow == 0) break;
            }
        }
        vis[u] = 0;
        return sum;
    }
    p64 work(int s,int t){
        p64 ans = {0,0};
        SPFA(s,t);
        while (dijkstra(s,t)){
            std::fill(vis.begin(),vis.end(),0);
            for (int i = 0;i <= n;i++){
                h[i] += dis[i];
                now[i] = head[i];
            }
            ans[0] += dfs(s,t,INF,ans[1]);
        }
        return ans;
    }
};
\end{lstlisting}

\subsubsection{基于 HLPP 的费用流（费用缩放）}

先用 \texttt{HLPP} 求出最大流把流值定住，再在残量网络上做费用缩放：势能与 $\varepsilon$ 都按 $n+1$ 倍整数存储，$\varepsilon$ 逐轮减半，每轮先把简化费用为负的残量弧饱和、再把产生的超额沿简化费用为负的弧压回平衡，$\varepsilon$ 降到 $1/(n+1)$ 时残量网络无负环，即得该流值下的最小费用，基于 \texttt{HLPP}。复杂度上界 $\mathcal{O}(nm\log (nC))$，$C$ 为边费用的最大绝对值；迭代轮数只与 $n$、$C$ 有关而与流量无关，流量大时远快于上面两种按流量逐次增广的 SSP，负费用乃至负环费用也能处理。

\texttt{work} 返回 \verb|{flow, cost}|，与 \texttt{SSP} 一致。费用绝对值大且点数多时注意 \texttt{long long}（量级 $|w|\cdot n^2$）。

\begin{lstlisting}
    using i64 = long long;
    using p64 = std::array<i64,2>;
    constexpr i64 INF = 1e18;
    struct MinCostFlow {
        int n,s = 0,t = 0;
        int maxh = 0,maxgaph = 0,workcnt = 0;
        struct Edge {
            int to = 0,nxt = 0;
            i64 f = 0,c = 0,w = 0;
        };
        std::vector<int> head,cur,h,ovList,ovNxt,gap,gapPrv,gapNxt;
        std::vector<i64> ov,pot;
        std::vector<bool> inq;
        std::vector<Edge> edge;
        std::queue<int> ovq;
    
        MinCostFlow(int _n):n(_n),head(n+1),cur(n+1),h(n+1),ovList(n+1,-1),
            ovNxt(n+1,-1),gap(n+1,-1),gapPrv(n+1,-1),gapNxt(n+1,-1),ov(n+1),
            pot(n+1),inq(n+1),edge(2) {}
    
        void connect(int u,int v,i64 c,i64 w){
            edge.push_back({v,head[u],c,c,w});
            head[u] = edge.size() - 1;
            edge.push_back({u,head[v],0,0,-w});
            head[v] = edge.size() - 1;
        }
    
        p64 work(int s_,int t_){
            s = s_,t = t_;
            i64 flow = maxFlow();
            ov[s] += flow;
            ov[t] -= flow;
            std::fill(pot.begin(),pot.end(),0);
            for (int u = 1;u <= n;++u){
                cur[u] = head[u];
                if (ov[u] > 0){
                    inq[u] = 1;
                    ovq.push(u);
                }
            }
            i64 mx = 0,eps = 1;
            for (int e = 2;e < (int)edge.size();++e) mx = std::max(mx,std::abs(edge[e].w));
            while (eps < mx * (n + 1)) eps <<= 1;
            for (;eps >= 1;eps >>= 1) refine(eps);
            i64 cost = 0;
            for (int e = 2;e < (int)edge.size();e += 2) cost += (edge[e].c - edge[e].f) * edge[e].w;
            return {flow,cost};
        }
    private:
        void setHeight(int x,int newh){
            if (~gapPrv[x]){
                if (gapPrv[x] == x){
                    gap[h[x]] = gapNxt[x];
                    if (~gapNxt[x]) gapPrv[gapNxt[x]] = gapNxt[x];
                } else {
                    gapNxt[gapPrv[x]] = gapNxt[x];
                    if (~gapNxt[x]) gapPrv[gapNxt[x]] = gapPrv[x];
                }
                gapPrv[x] = gapNxt[x] = -1;
            }
            h[x] = newh;
            if (h[x] >= n) return;
            maxgaph = std::max(maxgaph,h[x]);
            if (ov[x] > 0){
                maxh = std::max(maxh,h[x]);
                ovNxt[x] = ovList[h[x]];
                ovList[h[x]] = x;
            }
            gapNxt[x] = gap[h[x]];
            if (~gapNxt[x]) gapPrv[gapNxt[x]] = x;
            gap[h[x]] = gapPrv[x] = x;
        }
        void globalRelabel(){
            workcnt = maxh = maxgaph = 0;
            std::fill(h.begin(),h.end(),n);
            h[t] = 0;
            std::fill(gapPrv.begin(),gapPrv.end(),-1);
            std::fill(gapNxt.begin(),gapNxt.end(),-1);
            std::fill(gap.begin(),gap.end(),-1);
            std::fill(ovList.begin(),ovList.end(),-1);
            for (int u = 1;u <= n;++u) cur[u] = head[u];
            std::queue<int> que;
            que.push(t);
            while (!que.empty()){
                int x = que.front();
                que.pop();
                for (int e = head[x];e;e = edge[e].nxt){
                    int y = edge[e].to;
                    if (h[y] == n && y != s && edge[e ^ 1].f > 0){
                        setHeight(y,h[x] + 1);
                        que.push(y);
                    }
                }
            }
        }
        void pushFlow(int from,int e,i64 f){
            int v = edge[e].to;
            if (!ov[v] && v != t && h[v] < n){
                maxh = std::max(maxh,h[v]);
                ovNxt[v] = ovList[h[v]];
                ovList[h[v]] = v;
            }
            edge[e].f -= f;
            edge[e ^ 1].f += f;
            ov[from] -= f;
            ov[v] += f;
        }
        void discharge(int x){
            int nh = n,cur0 = cur[x];
            for (int e = cur0;e;e = edge[e].nxt){
                if (edge[e].f > 0){
                    int y = edge[e].to;
                    if (h[x] == h[y] + 1){
                        pushFlow(x,e,std::min(ov[x],edge[e].f));
                        if (ov[x] == 0){
                            cur[x] = e;
                            return;
                        }
                    } else {
                        nh = std::min(nh,h[y] + 1);
                    }
                }
            }
            for (int e = head[x];e != cur0;e = edge[e].nxt){
                if (edge[e].f > 0) nh = std::min(nh,h[edge[e].to] + 1);
            }
            cur[x] = head[x];
            ++workcnt;
            int oldh = h[x],head0 = gap[oldh];
            if (~gapNxt[head0]){
                setHeight(x,nh);
            } else {
                for (int i = oldh;i <= maxgaph;++i){
                    for (int j = gap[i];~j;){
                        int nxt = gapNxt[j];
                        h[j] = n;
                        gapPrv[j] = gapNxt[j] = -1;
                        j = nxt;
                    }
                    gap[i] = -1;
                }
                maxgaph = oldh - 1;
            }
        }
        i64 maxFlow(){
            globalRelabel();
            for (int e = head[s];e;e = edge[e].nxt){
                if (edge[e].f) pushFlow(s,e,edge[e].f);
            }
            for (;maxh >= 0;--maxh){
                while (~ovList[maxh]){
                    int x = ovList[maxh];
                    ovList[maxh] = ovNxt[x];
                    discharge(x);
                    if (workcnt > (n << 2)) globalRelabel();
                }
            }
            return ov[t];
        }
        i64 red(int u,int e){
            return (n + 1) * edge[e].w + pot[u] - pot[edge[e].to];
        }
        void pushCost(int u,int e,i64 f){
            int v = edge[e].to;
            edge[e].f -= f;
            edge[e ^ 1].f += f;
            ov[u] -= f;
            ov[v] += f;
            if (ov[v] > 0 && !inq[v]){
                inq[v] = 1;
                ovq.push(v);
            }
        }
        void dischargeCost(int u,i64 eps){
            while (ov[u] > 0){
                if (!cur[u]){
                    i64 mn = INF;
                    for (int e = head[u];e;e = edge[e].nxt){
                        if (edge[e].f > 0) mn = std::min(mn,red(u,e));
                    }
                    pot[u] -= mn + eps;
                    cur[u] = head[u];
                }
                int e = cur[u];
                if (edge[e].f > 0 && red(u,e) < 0) pushCost(u,e,std::min(ov[u],edge[e].f));
                else cur[u] = edge[e].nxt;
            }
        }
        void refine(i64 eps){
            for (int u = 1;u <= n;++u){
                for (int e = head[u];e;e = edge[e].nxt){
                    if (edge[e].f > 0 && red(u,e) < 0) pushCost(u,e,edge[e].f);
                }
            }
            while (!ovq.empty()){
                int u = ovq.front();
                ovq.pop();
                inq[u] = 0;
                dischargeCost(u,eps);
            }
        }
    };
\end{lstlisting}

